% #1114 复现：命令正文里用户写的盒子、脚注被当作命令自己的首尾类别
% 每组上一行是命令包裹的写法，下一行是删去外层命令的直接输入；
% 正确时两行等宽。上一行用灰框标出盒子（只为显示位置，宽度按不加框的写法量取），
% 右侧给出两行的宽度差。
\documentclass{article}
\usepackage{xeCJK}
\usepackage{xcolor}
\usepackage[colorlinks,urlcolor=blue]{hyperref}
\setCJKmainfont{FandolSong-Regular.otf}
\pagestyle{empty}
\setlength{\parindent}{0pt}
\newcommand{\HB}{\fboxsep=0pt\fboxrule=0.2pt\textcolor{gray}{\fbox{\textcolor{black}{y}}}}
\newcommand{\RawHB}{\hbox{y}}
\newcommand{\PB}{\parbox{2.5cm}{y}}
\newcommand{\FN}{\footnote{y}}
\renewcommand{\thefootnote}{1}
\newcommand{\SH}{\setbox0=\hbox{y}}
\newcommand{\SBX}{\sbox0{y}}
\newlength{\candwd}
\newlength{\directwd}
% #1 名称，#2 包裹写法（宽度量取），#3 直接输入，#4 可视化用写法
\newcommand{\Row}[4]{%
  \settowidth{\candwd}{#2}\settowidth{\directwd}{#3}%
  \par\medskip
  \makebox[11em][l]{\small\texttt{#1}}%
  \makebox[9em][l]{\large #4}%
  \makebox[9em][r]{\small\the\candwd}%
  \ifdim\candwd=\directwd \textcolor{green!50!black}{\quad 一致}%
  \else \textcolor{red}{\quad 差 \the\dimexpr\candwd-\directwd\relax}\fi
  \typeout{ROW #1: cand=\the\candwd, direct=\the\directwd}%
  \par
  \makebox[11em][l]{}%
  \makebox[9em][l]{\large #3}%
  \makebox[9em][r]{\small\the\directwd}\par}
\begin{document}
\section*{xeCJK \#1114 复现（master \texttt{e07fcbf2}）}
\Row{href + \textbackslash hbox}{中\href{u}{\RawHB}文}{中\RawHB{}文}{中\href{u}{\HB}文}
\Row{href + \textbackslash hbox（西文侧）}{x\href{u}{\RawHB}中}{x\RawHB{}中}{x\href{u}{\HB}中}
\Row{href + \textbackslash parbox}{中\href{u}{\PB}文}{中\PB{}文}{中\href{u}{\PB}文}
\Row{mbox + \textbackslash hbox}{中\mbox{\RawHB}文}{中\RawHB{}文}{中\mbox{\HB}文}
\Row{href + 脚注}{x\href{u}{\FN}中}{x\FN{}中}{x\href{u}{\FN}中}
\Row{href + \textbackslash setbox}{x\href{u}{\SH 文}x}{x\SH 文x}{x\href{u}{\SH 文}x}
\Row{href + \textbackslash sbox（对照）}{x\href{u}{\SBX 文}x}{x\SBX 文x}{x\href{u}{\SBX 文}x}
\end{document}
